% FONTE TEMA https://github.com/matze/mtheme
\documentclass[aspectratio=1610]{beamer}
%\documentclass[aspectratio=1610, handout]{beamer}
\usepackage[utf8]{inputenc}
\usepackage{ragged2e}
\usepackage{xcolor}
\usepackage[italian]{babel}
\usepackage{multirow}
\usepackage{silence}
\usepackage{tikz}
\usepackage[normalem]{ulem}
\WarningFilter{beamer}{}
\WarningFilter{metropolis}{}
\usetheme[progressbar=frametitle,titleformat=smallcaps]{metropolis}
\setbeamertemplate{frame numbering}[fraction]
\setbeamercovered{dynamic}
\definecolor{rosso}{RGB}{255, 0, 0}
\definecolor{giallo}{RGB}{254,212,23}
\hypersetup{colorlinks=true,linkcolor=black,urlcolor=rosso}
\setbeamercolor{palette primary}{fg=black, bg=giallo}
\setbeamercolor{background canvas}{bg=white}
\setbeamercolor{normal text}{fg=black}
\setbeamercolor{progress bar}{fg=rosso}
\setbeamercolor{framesubtitle}{fg=rosso}
\setbeamercolor{normal text .dimmed}{fg=giallo}
\setbeamercolor{block title alerted}{fg=rosso, bg=giallo}
\setbeamerfont{caption}{size=\tiny}
\setbeamerfont{caption name}{size=\tiny}
\setlength{\abovecaptionskip}{0pt}
\makeatletter
\metroset{block=fill}
\setlength{\metropolis@progressinheadfoot@linewidth}{1pt} 
\setlength{\metropolis@progressonsectionpage@linewidth}{1pt}
\setlength{\metropolis@titleseparator@linewidth}{1pt}
\makeatother

\title{SEDUTO E MUTO}
\subtitle{Gioco in solitaria}
\date{}
\institute{Materiale necessario: un foglio e una penna, un conto alla rovescia}

\begin{document}

\begin{frame}[plain, noframenumbering]
    \titlepage
\end{frame}

\begin{frame}{SEDUTO E MUTO - OBIETTIVO DEL GIOCO}
    \begin{minipage}{0.98\linewidth}
        \centering
        \huge
        \textit{Ogni giocatore deve riuscire entro il tempo limite di 2 minuti a far sedere più compagni 
        di classe possibile, cercando di lasciare in piedi un solo compagno target.}\\
    \end{minipage}\\
\end{frame}

\begin{frame}{SEDUTO E MUTO - REGOLE DEL GIOCO}
    GIOCATORE:
    \begin{enumerate}
        \item Il giocatore va alla cattedra e scrive sul foglio il nome del compagno target da far 
        rimanere in piedi (non può essere un target già scelto da un altro giocatore);
        \item Il giocatore avvia il conto alla rovescia di 2 minuti e cerca di far sedere più compagni 
        possibile ponendo loro delle domande (S\'I/NO);
        \item Le domande non possono essere troppo specifiche (approvate dal Prof.);
        \item Se il compagno target si siede, il giocatore finisce il turno;
        \item Se tutti i compagni si siedono, il giocatore finisce il turno;
        \item Se rimane in piedi solo il compagno target, il giocatore finisce il turno.
    \end{enumerate}
\end{frame}

\begin{frame}{SEDUTO E MUTO - REGOLE DEL GIOCO}
    CLASSE:
    \begin{enumerate}
        \item Inizialmente i giocatori devono rimanere in piedi;
        \item Se la risposta alla domanda del giocatore alla cattedra è affermativa, il giocatore deve sedersi;
        \item Se la risposta alla domanda del giocatore alla cattedra è negativa, il giocatore deve rimanere in piedi;
        \item Terminato il turno, il giocatore target prende il posto del giocatore alla cattedra e inizia il suo turno.
    \end{enumerate}
\end{frame}

\begin{frame}{SEDUTO E MUTO - PUNTEGGI}
    \begin{itemize}
        \item Per entrare in classifica bisogna riuscire a fare almeno un punto;
        \item Ogni compagno che rimane in piedi alla fine del turno vale 1 punto;
        \item Vince chi ottiene il punteggio più basso.
    \end{itemize}
\end{frame}
\end{document}